\documentclass{article}
\usepackage[utf8]{inputenc}
\usepackage[T1]{fontenc}
\usepackage[margin=2cm]{geometry}
\usepackage{hyperref}
\sloppy
\begin{document}
\section*{Historical source}
Encyclopaedia Britannica scan, source row 1904. Parse 3.5 extracted the text. Super 3.5 VL corrected it from an enlarged crop of each block. It can still contain recognition errors.
Page headers, footers, and page numbers are removed.

\subsection*{Untitled continuation (E0166)}
\textit{Continues an entry from an earlier page.}

LYR [ 4329 ] LYR

which he used to frequent, at the foot of a mountain in Arcadia. Just as he was going in, he found a tortoise feeding at the entrance of his cave: he killed the poor creature, and, perhaps, eat the flesh of it. As he was diverting himself with the shell, he was mightily pleased with the noise it gave from its concave figure. He had possibly been cunning enough to find out, that a thong pulled tight and fastened at each end, when struck by the finger, made a sort of musical found. However that was, he went immediately to work, and cut several things out of the hides he had lately stolen, and fastened them as tight as he could to the shell of this tortoise; and, in playing with them, made a new kind of music with them to divert himself in his retreat. This, considered only as an account of the first invention of the lyre, is not altogether unnatural.

The most ancient representations of this instrument agree very well with the account of its invention: the lyre, in particular on the old celestial globes, was represented as made of one entire shell of a tortoise; and that Amphion in the celebrated group of the Dirce, or Toro, in the Farnese palace at Rome, which is of Greek sculpture, and very high antiquity, is figured in the same manner. See Plate CLXI. fig. 9.

There have, however, been many other claimants to the seven-stringed lyre. For though Mercury invented this instrument in the manner already related, it is said he afterwards gave it to Apollo, who was the first that played upon it with method, and made it the constant companion of poetry. According to Homer's account of this transaction, in his hymn to Mercury, it was given by that god to Apollo, as a peace offering and indemnification for the oxen which he had stolen from him:

To Phoebus Maia's son presents the lyre,
A gift intended to appease his ire.

A gift intended to appease his ire.
The god receives it gladly, and displays
The novel instrument a thousand ways:
With dexterous skill the plectrum wields; and
With voice accordant to the trembling strings,
Such strains as gods and men approved, from w
The sweet alliance sprung of found and fence.

Diodorus informs us, that Apollo soon repenting of the cruelty with which he had treated Marsyas in consequence of their musical contest, broke the strings of the lyre, and by that means put a stop for a time to any further progress in the practice of that new instrument. "The muses (adds he) afterwards added to this instrument the string called mese; Linus, that of lichanos; and Orpheus and Thamyras, those strings which are named hypate and parhypate (a).

Again, many ancient and reputable authors tell

us, that, before the time of Terpander, the Grecian lyre had only four strings; and, if we may believe Suidas, it remained in this state 856 years, from the time of Amphion, till Terpander added to it three new strings, which extended the musical scale to a heptachord, or seventh, and supplied the player with two conjoint tetrachords. It was about 150 years after this period, that Pythagoras is said to have added an eighth string to the lyre, in order to complete the octave, which consisted of two disjunct tetrachords.

Bochthia gives a different history of the feels, and tells me, that the flyton did no long remains in fact narrow limits as a tetrachold. Chorobus, the fan of Athis, or Arys, king of Lydia, added a fifth firing; Hyguns, a firstly, Terpandera, a favorite and, at large), Erythema of Samoa, an eighth. But all their accounts are irreconcileable with Momars, hymn to Mercury, where the chelys, or reduce, the invention of which he earlier or than god, is fast to have had been fittings. There are many clamans among the musicians of ancient Greece, to the things that were afterwards added to thefly, by which the date, in the time of Alfoxcoras, was extended to two offices. Athamus, more than once, peaks of the nine fingerling inflamments, and Lon OlChoja rtagies and lyrics pace and philosopher, who fitf recited his pieces in the 8zd olympiad, 452 B. C. mentions, in some refer quoted in Dr. David, the ten- firginged lyrs, a prodt that, the third conjoint tetr- chord was added to the scale in his time, which was about 50 years after Pythagoras is fuppelled to have confinational the death.

The different claimants among the Greeks to the fame musical discoveries, only prove, that music was cultivated in different countries; and that the inhabitants of each country invented and improved their own instruments, some of which happening to resemble those of other parts of Greece, rendered it difficult for historians to avoid attributing the fame invention to different persons. Thus the single flute was given to Minerva, and to Marsyas; the syrinx, or fistula, to Pan, and to Cybele; and the lyre, or cithara, to Mercury, Apollo, Amphion, Linus, and Orpheus. Indeed, the mere addition of a string or two to an instrument without a neck, was so obvious and easy, that it is scarce possible not to conceive many people to have done it at the same time.

With respect to the form of the ancient lyre; as little agreement is to be found among authors as about the number of strings. The best evidences concerning it are the representations of that instrument in the hands of ancient statues, bas-reliefs, \&c. See Plate CLXI.

(a) It has been already related, that the lyre invented by the Egyptian Mercury had but three strings; and by putting these two circumstances together, Dr Burney observes, we may perhaps acquire some knowledge of the progress of music, or, at least, of the extension of its scale, in the highest antiquity.

Mese, in the Greek music, is the fourth found of the second tetrachord of the great system, and first tetrachord invented by the ancients, answering to our A, on the fifth line in the base. If this found then was added to the former three, it proves two important points: first, that the most ancient tetrachord was that from E in the base to A; and that the three original strings in the Mercurian and Apollonian lyre were tuned E, F, G, which the Greeks called Hypate Meson, Parhypate Meson, Meson Diatonos. The addition therefore of Mese to these, completed the first and most ancient tetrachord, E, F, G, A.

The firing lichanos, then, being added to these, and answering to our D on the third line in the base, extended the compass downwards, and gave the ancient lyre a regular series of five sounds in the Dorian mode, the most ancient of all the Greek modes; and the two frings called Hypate and Parhypate, corresponding with our B and C in the base, completed the heptachord, or seven sounds, B, C, D, E, F, G, A, a compass that received no addition till after the time of Pindar, who calls the instrument then in use the seven-tongued lyre.


\section*{Source record}
\begin{description}
\item[Scan ID] \texttt{britannica:309f765b1ecc83e0a7f7762e44b4a927a7555443:train:1904}
\item[Image file] \texttt{page-1904-759854fd95bb.jpg}
\item[Image SHA-256] \texttt{759854fd95bb63859232323bbddae58dd20122c44ba36e429b4243846268277d}
\item[Dataset] NationalLibraryOfScotland/encyclopaedia\_britannica\_illustrated at revision \texttt{309f765b1ecc83e0a7f7762e44b4a927a7555443} (CC0-1.0)
\item[Extraction] model \texttt{nvidia/nemotron-3.5-parse}, raw reply \texttt{extract/raw/page-1904-759854fd95bb.json}
\item[Raw reply SHA-256] \texttt{2d40c0bfd0dd16197cef52e50485c3e171e9057e89f36999a20f2b221d74c3f6}
\item[Correction] model \texttt{nvidia/Nemotron-3.5-Super-VL}, block crops, raw replies \texttt{correct/raw/crop\_nothink/page-1904-759854fd95bb-b*.json}; super\_vl: 15
\item[Running header] none
\end{description}
\end{document}
